Write \(b_t=B_{\cdot,t}\), \(h_t=h_e(t)\), and form the finite observable--target hull
\[
\mathcal C=\{(Bv,h_e^{\top}v):v\in\Delta_{\mathrm{SLCC}}\}
=\operatorname{conv}\{(b_t,h_t):t\in\mathcal T_{\mathrm{SLCC}}\}.
\tag{1}
\]
The simplex is compact and the map in (1) is linear.  Therefore \(\mathcal C\), every fiber \(\{v\in\Delta_{\mathrm{SLCC}}:Bv=p\}\), and its image under \(v\mapsto h_e^{\top}v\) are compact.  For \(p\in\mathfrak O\), the last image is also nonempty and convex, hence a compact interval.  Its endpoints are attained and equal \(L(p)\) and \(U(p)\) by \(\text{def:sharp-endpoint-programs}\).

We derive duality without silently normalizing the mass multiplier.  The membership condition \((x,z)\in\mathcal C\) is the finite linear system
\[
v_t\geq0\quad(t\in\mathcal T_{\mathrm{SLCC}}),
\qquad
\mathbf 1^{\top}v=1,
\qquad
Bv=x,
\qquad
h_e^{\top}v=z.
\tag{2}
\]
Replace equalities by pairs of inequalities and eliminate the coordinates of \(v\) successively.  At each elimination step, inequalities with positive coefficient give upper bounds, those with negative coefficient give lower bounds, and those with zero coefficient remain; existence of the eliminated coordinate is equivalent to every retained lower bound being at most every retained upper bound.  Thus each step preserves equivalence and finiteness.  After finitely many steps, there exist finite \(R\), vectors \(a_r\), and scalars \(d_r,c_r\) such that
\[
\mathcal C=\{(x,z):a_r^{\top}x+d_rz\leq c_r\text{ for every }r\in R\}.
\tag{3}
\]
The valid inequalities \(-z\leq0\) and \(z\leq1\), which follow from \(h_t\in\{0,1\}\), may be added to (3).  Hence \(R_-:=\{r:d_r<0\}\) and \(R_+:=\{r:d_r>0\}\) are nonempty.  For fixed \(p\in\mathfrak O\), the constraints with \(d_r=0\) are already satisfied, and solving (3) for \(z\) gives
\[
L(p)=\max_{r\in R_-}\frac{c_r-a_r^{\top}p}{d_r},
\qquad
U(p)=\min_{r\in R_+}\frac{c_r-a_r^{\top}p}{d_r}.
\tag{4}
\]

For \(r\in R_-\), write its affine function in (4) as \(\alpha_r+\lambda_r^{\top}p\).  Substituting every generator \((b_t,h_t)\) into the defining inequality in (3), and dividing by \(d_r<0\), gives
\[
\alpha_r+\lambda_r^{\top}b_t\leq h_t
\qquad\text{for every }t.
\tag{5}
\]
Thus \((\alpha_r,\lambda_r)\in\mathcal D_L\).  Conversely, if \((\alpha,\lambda)\in\mathcal D_L\) and \(Bv=p\), then
\[
\alpha+\lambda^{\top}p
=\sum_t v_t(\alpha+\lambda^{\top}b_t)
\leq\sum_tv_th_t=h_e^{\top}v.
\tag{6}
\]
Taking the minimum in (6) and then choosing an index attaining the finite maximum in (4) proves
\[
L(p)=\max_{(\alpha,\lambda)\in\mathcal D_L}
\{\alpha+\lambda^{\top}p\},
\tag{7}
\]
and exhibits at least one optimizer.  Repeating (5)--(7) for \(r\in R_+\), with inequalities reversed after division by \(d_r>0\), gives
\[
U(p)=\min_{(\alpha,\lambda)\in\mathcal D_U}
\{\alpha+\lambda^{\top}p\},
\tag{8}
\]
again with attainment.

Equation (4) is a maximum or minimum of finitely many affine functions.  Therefore \(L\) and \(U\) are finite piecewise-affine and continuous on all of their relative domain \(\mathfrak O\), including its boundary.  Applying exactly the same argument to the finite matrix \(K_{\mathrm{blind}}B\) proves attainment, finite-polyhedral duality, and continuity for \(L_{\mathrm{blind}}\) and \(U_{\mathrm{blind}}\) on \(K_{\mathrm{blind}}\mathfrak O\).